\clearpage
\begingroup
\fancyhead[L]{\small BEMERKUNGEN}
\renewcommand{\headrulewidth}{0pt}
\SourceLeaf{607}
\needspace{6\baselineskip}
\pdfbookmark[1]{Bemerkungen zu den Störungen der Pallas}{pallas-bemerkungen}
\begin{center}\large BEMERKUNGEN.\end{center}
\begingroup
\newcommand{\BIntroJ}{\text{\mdseries\AstroSymbols ♃}}
\newcommand{\BIntroP}{\text{\mdseries\AstroSymbols ⚴}}

\par\noindent\hspace*{\parindent}Von \textsc{Gauss}’ Untersuchungen über die Störungen der Pallas finden sich im Nachlass fast ausschliesslich
numerische Rechnungen ohne erklärenden Text; diese stehen auf einer grossen Menge einzelner Zettel,
welche im Archiv fünf Kapseln füllen; die Berechnung der Oppositionen und die Vergleichung mit den
Beobachtungen hat \textsc{Gauss} indessen auch vielfach in seinen Handbüchern ausgeführt. Die grosse Menge des
Materials erschwerte die Sichtung wesentlich; doch war es schliesslich möglich, namentlich durch gleichzeitiges
Studium des Briefwechsels, das Ganze in zusammenhängende Form zu bringen.

\par\noindent\hspace*{\parindent}Im Vorstehenden sind die Stellen aus dem Briefwechsel (S.~413—438) vorangestellt, welche den
historischen Verlauf der \textsc{Gauss}schen Untersuchungen klarlegen; über ihren Zusammenhang mit dem später
abgedruckten Text sei Folgendes bemerkt:

\medskip
\par\noindent\hspace*{\parindent}\textsc{Gauss} hatte während der Berechnung der Ceresstörungen den Plan gefasst, eine Tafel herzustellen,
welche die Entwickelung der Störungsfunction für die gesamten bis dahin entdeckten kleinen Planeten
erleichtern sollte; er schreibt darüber an \textsc{Olbers} am 3.~September, 7.~und 29.~October 1805. Auch findet
sich im Handbuch B\,e eine Reihe von Aufzeichnungen über die Methode, die \textsc{Gauss} hier anzuwenden gedachte.
Die Tafel, bei deren Berechnung \textsc{Bessel} half, ist auch vollendet worden und im Nachlass vorhanden.
Da \textsc{Gauss} dann aber bei der Berechnung der Pallasstörungen andere Methoden anwandte, so ist
sie niemals zur Anwendung gekommen, und aus diesem Grunde wurde auch auf ihren Abdruck verzichtet.

\par\noindent\hspace*{\parindent}Die Briefe an \textsc{Olbers} von Januar 1806 bis October 1810 beziehen sich auf Vorarbeiten zur Berechnung
der Störungen und zwar theils auf Untersuchungen über die Genauigkeit, mit der sich die Beobachtungen
bei Anwendung der Methode der kleinsten Quadrate darstellen lassen, theils auf die analytischen
Grundlagen zur Entwickelung der Störungsfunction. Ende October 1810 hat \textsc{Gauss} dann die erste Rechnung
der speciellen Störungen der Pallas (abgedruckt S.~473) begonnen, auf welche sich die Briefe vom
November 1810 beziehen. Der Brief an \textsc{Olbers} vom 13.~December 1810 betrifft die zweite Rechnung der
speciellen Störungen (S.~483), die von August und October 1811 betreffen die erste Rechnung der allgemeinen
Störungen (S.~489). Im April 1812 begann \textsc{Gauss} die zweite Rechnung der allgemeinen Störungen,
die Juli 1813 vollendet wurde. Aus der Zeit von März 1812 bis Januar 1813 haben ausser dem Brief an
\textsc{Bessel} vom 5.~Mai 1812 nur die Briefe von \textsc{Olbers} an \textsc{Gauss} abgedruckt werden können, da die entsprechenden
von \textsc{Gauss} an \textsc{Olbers} verloren gegangen sind; unter den verlorenen befindet sich auch der,
in welchem \textsc{Gauss}, wie im genannten Brief an \textsc{Bessel}, Mittheilung von seiner Entdeckung der Rationalität
der Umlaufszeiten von Pallas und Jupiter macht. \textsc{Gauss} hat diese Entdeckung zuerst also vor der
Ausführung der zweiten Rechnung der allgemeinen Störungen im März oder Anfang April 1812 gemacht,
als er mit dem kleinen Divisor \(18\BIntroJ-7\BIntroP\) noch garnicht gerechnet hatte; es bleibt unklar (vgl. aber das
weiter unten S.~604 Gesagte), in welcher Weise er die Rationalität im speciellen gefunden hat; die Notizen,
\par\smallskip\noindent VII.\hfill 76\par
\noindent\SourceLeaf{608}die sich hierüber im Nachlass fanden und die S.~558 abgedruckt sind, stammen aus späterer Zeit, jedenfalls
\textls[120]{nach} der Ausführung mindestens eines grossen Theils der zweiten Rechnung.

\par\noindent\hspace*{\parindent}Besondere Sorgfalt verwandte \textsc{Gauss} auf die Vergleichung seiner Resultate mit den Beobachtungen,
resp. auf die Bestimmung der Normalelemente, welche in Verbindung mit den gerechneten Störungen die
Beobachtungen möglichst gut darstellten. Man ist somit in der Lage, sich ein sehr gutes Bild darüber
zu machen, wie genau die Beobachtungen mit der Theorie übereinstimmen; die Differenzen zwischen beiden
hat \textsc{Gauss} meist schon selbst publicirt (s.~Band VI); ausserdem findet man sie, ausser bei den Vergleichungen
mit den Oppositionen, auch in den oben abgedruckten Briefstellen. Ausser den im Anschluss an die Bestimmung
der Störungen oben abgedruckten Vergleichungen mit den Oppositionen findet sich im Nachlass
eine sehr grosse Menge solcher Vergleichungen, auch für spätere Oppositionen, theils mit Berücksichtigung
der Saturnstörungen theils ohne diese, theils mit der alten Jupitermasse, theils mit verbesserter.
Am vollständigsten ist eine Vergleichung mit den ersten 10 resp. 12 Oppositionen (1802 bis 1817) im Handbuch
Bf erhalten; sie ist nicht abgedruckt worden, weil sie in der Methode nichts Neues bietet und ihre
Resultate aus dem S.~430 abgedruckten Briefe an \textsc{Olbers} hervorgehen. Diese Vergleichungen mit den Oppositionen
hat \textsc{Gauss} sowohl an seine speciellen Störungen angeschlossen, die er noch bis zum Jahre 1816
fortgesetzt zu haben scheint, als auch an die allgemeinen. Bei der Berechnung der Störungsbeträge aus
den Ausdrücken für die allgemeinen Störungen haben ihm \textsc{Westphal}, \textsc{Encke} und \textsc{Nicolai} wesentliche
Hülfe geleistet; hiezu wurden, so lange die Tafeln (S.~572) noch nicht fertig waren, Gruppen von Störungsgleichungen,
je nach ihrer Grösse, auf Zettelchen zusammengefasst und »Manipel« genannt. Jeder Manipel
(mit Ausnahme des letzten) enthielt zehn Störungsglieder; so lauten z.~B. die ersten Manipel für die
Länge des Perihels folgendermassen:

\needspace{19\baselineskip}
\begin{center}\textit{Länge des Perihels.}\end{center}
\begingroup\small
\setlength{\arraycolsep}{3pt}
\begin{center}\textit{Manipel 1.}\end{center}
\[
\begin{array}{ccc}
\left[\begin{gathered}\text{\itshape Nr. der Stö-}\\\text{\itshape rungsgleichung}\end{gathered}\right]&
\text{\itshape [Argument]}&
\left[\begin{gathered}\text{\itshape Log. des}\\\text{\itshape Coefficienten}\end{gathered}\right]\\
\mathit{27}&\mathit{2}\BIntroJ-\BIntroP+\mathit{236^0\,43'\,49''}&\mathit{2{,}88715}\\
\mathit{66}&\mathit{5}\BIntroJ-\mathit{2}\BIntroP+\mathit{303\ \,26\ \,18}&\mathit{2{,}81492}\\
\mathit{40}&\mathit{3}\BIntroJ-\BIntroP+\mathit{94\ \,47\ \,46}&\mathit{2{,}72535}\\
\mathit{52}&\mathit{4}\BIntroJ-\mathit{2}\BIntroP+\mathit{288\ \,15\ \,4}&\mathit{2{,}16040}\\
\mathit{39}&\mathit{3}\BIntroJ-\mathit{2}\BIntroP+\mathit{260\ \,27\ \,22}&\mathit{1{,}95730}\\
\mathit{26}&\mathit{2}\BIntroJ-\mathit{2}\BIntroP+\mathit{342\ \,9\ \,24}&\mathit{1{,}94234}\\
\mathit{28}&\mathit{2}\BIntroJ\hspace{2em}+\mathit{101\ \,42\ \,47}&\mathit{1{,}91320}\\
\mathit{100}&\mathit{8}\BIntroJ-\mathit{3}\BIntroP+\mathit{168\ \,38\ \,33}&\mathit{1{,}90406}\\
\mathit{53}&\mathit{4}\BIntroJ-\BIntroP+\mathit{106\ \,12\ \,53}&\mathit{1{,}88715}\\
\mathit{14}&\BIntroJ\hspace{2em}+\mathit{291\ \,57\ \,59}&\mathit{1{,}85959}
\end{array}
\]
\begin{center}\textit{Manipel 2.}\end{center}
\[
\begin{array}{ccc}
\left[\begin{gathered}\text{\itshape Nr. der Stö-}\\\text{\itshape rungsgleichung}\end{gathered}\right]&
\text{\itshape [Argument]}&
\left[\begin{gathered}\text{\itshape Log. des}\\\text{\itshape Coefficienten}\end{gathered}\right]\\
\mathit{78}&\mathit{6}\BIntroJ-\mathit{2}\BIntroP+\mathit{137^0\,59'\,40''}&\mathit{1{,}77929}\\
\mathit{13}&\BIntroJ-\BIntroP+\mathit{309\ \,34\ \,12}&\mathit{1{,}76702}\\
\mathit{1}&\hspace{1.5em}\BIntroP+\mathit{321\ \,28\ \,57}&\mathit{1{,}75857}\\
\mathit{89}&\mathit{7}\BIntroJ-\mathit{3}\BIntroP+\mathit{334\ \,36\ \,27}&\mathit{1{,}75507}\\
\mathit{38}&\mathit{3}\BIntroJ-\mathit{3}\BIntroP+\mathit{340\ \,39\ \,45}&\mathit{1{,}68732}\\
\mathit{77}&\mathit{6}\BIntroJ-\mathit{3}\BIntroP+\mathit{320\ \,40\ \,29}&\mathit{1{,}59337}\\
\mathit{65}&\mathit{5}\BIntroJ-\mathit{3}\BIntroP+\mathit{306\ \,34\ \,54}&\mathit{1{,}57175}\\
\mathit{122}&\mathit{10}\BIntroJ-\mathit{4}\BIntroP+\mathit{18\ \,48\ \,20}&\mathit{1{,}54904}\\
\mathit{41}&\mathit{3}\BIntroJ\hspace{2em}+\mathit{87\ \,0\ \,49}&\mathit{1{,}50939}\\
\mathit{12}&\BIntroJ-\mathit{2}\BIntroP+\mathit{242\ \,28\ \,11}.&\mathit{1{,}43467}
\end{array}
\]
\endgroup
\begin{center}\textit{u. s. w.}\end{center}

\par\noindent Für jeden solcher Manipel wurden die einzelnen Störungsglieder berechnet und dann die Beträge der
Manipel addirt, um die Gesammtstörungen eines jeden Elementes zu finden. Specielle Störungen für spätere
Jahre hat auch \textsc{Encke} gerechnet, s.~Berliner Astronomisches Jahrbuch für 1837—38.

\par\noindent\hspace*{\parindent}Die erste Untersuchung über die Jupitermasse ist S.~561 abgedruckt. Der Verlauf von \textsc{Gauss}’ Arbeiten
hierüber ergibt sich aus dem Briefwechsel S.~424—429; auffallend ist, wie genau \textsc{Gauss} den Werth
dieser Grösse zuletzt im Briefe an \textsc{Olbers} vom 24.~Juli 1816 auf \(\frac{1}{1030}\) angibt.

\par\noindent\hspace*{\parindent}Auf die S.~572 abgedruckten Tafeln der Störungen durch Jupiter beziehen sich die Briefe von
\SourceLeaf{609}Februar 1817 bis März 1818; der erste Plan einer solchen Tafel wird bereits im Briefe an \textsc{Olbers} vom
8.~April 1813 erwähnt, in welchem \textsc{Gauss} auch die Absicht ausspricht, die Störungen durch Saturn und
Mars berechnen zu lassen.

\par\noindent\hspace*{\parindent}Über die Ausarbeitung seiner Methode, als deren Beginn wir die »Exposition d’une nouvelle méthode
etc.« anzusehen haben, finden sich Andeutungen in den Briefen vom 12.~December 1813 und vom 25.~September
1814. Dass diese mit der Absicht der Bewerbung um den von der französischen Akademie ausgesetzten
Preis begonnen wurde, geht schon aus ihrer Abfassung in französischer Sprache hervor. Wie sich
dann die Hoffnungen der astronomischen Welt, der Preis möchte \textsc{Gauss} zufallen, durch lange Jahre hinzogen
und doch schliesslich enttäuscht wurden, zeigt der Briefwechsel deutlich genug.
\endgroup
\clearpage
\endgroup
